% Ethics report: Spotify genre classifier (ML in Production Environment)
% Ready for Overleaf: it compiles with pdfLaTeX (Overleaf's default), XeLaTeX and LuaLaTeX.
\documentclass[11pt,a4paper]{article}

\usepackage{iftex}
\ifPDFTeX
  \usepackage[T1]{fontenc}
  \usepackage{lmodern}
\else
  \usepackage{fontspec}
\fi
\usepackage[english]{babel}
\usepackage[a4paper,margin=2.2cm]{geometry}
\usepackage{setspace}
\setstretch{1.1}
\usepackage{microtype}
\usepackage{parskip}
\usepackage{csquotes}
\usepackage{enumitem}
\usepackage{booktabs}
\usepackage{float}
\usepackage[font=small,skip=8pt]{caption}
\usepackage{xurl}                    % long links can break over lines
\usepackage[hidelinks,pdftitle={Ethics report: Spotify genre classifier}]{hyperref}
\pdfstringdefDisableCommands{\def\enquote#1{"#1"}}
\urlstyle{same}                      % links in the normal font

\usepackage{titlesec}
\titleformat*{\section}{\large\bfseries}
\titlespacing*{\section}{0pt}{2.2ex plus 1ex minus .2ex}{0.6ex plus .2ex}
\setcounter{secnumdepth}{0}          % headings without numbers
\setlist[itemize]{itemsep=4pt,parsep=0pt,topsep=6pt}
\renewcommand{\arraystretch}{1.15}    % a bit more room in the tables
\setlength{\emergencystretch}{1.5em} % no lines that stick out into the margin

\begin{document}

\begin{center}
  {\Large\bfseries Ethics report: Spotify genre classifier\par}
  \medskip
  ML in Production Environment\par
\end{center}
\medskip

\section{Introduction}

The app guesses the genre of a song from how it sounds. Behind it there is an XGBoost model, a Flask backend with a REST API and a React frontend, and with it you can make predictions, look at the test results and the feature importance, and change the thresholds. The lecture makes the point that the labels, the metrics and the thresholds in a machine learning project are never completely neutral choices. This report goes through them using the four steps from the lecture, map, measure, manage and govern, and most of the focus is on the thresholds.

\section{Who is affected?}

Music recommendation is at the low end of the lecture's stake spectrum, where the main worries are privacy and filter bubbles, and a genre classifier that only gives listeners a hint fits in the same place. Under the EU AI Act it would most likely count as minimal risk. Things change if a streaming service used the same model to decide which playlists a song ends up on, because then a wrong label would cost the artist listeners and money. The people affected are listeners and artists, and the worst thing that could realistically happen is that the same groups of artists get the wrong label over and over again. That is why the model works as a hint, but not as a judge, since genre is partly a question of culture and taste.

\section{Where things can go wrong}

The lecture lists five places where harm can get into a machine learning pipeline, and all five show up in this project.

The first is the data. The model only knows 6 of the 114 genres in the Kaggle dataset, all of them Western, which leaves 4,650 songs. Just like in the Gender Shades study, where face classifiers worked worse for women with darker skin, one overall number hides how different groups are doing. The app gets 80\% of the test songs right, and almost all of the Western classical songs (121 of 123), but only 23 of the 38 Indian songs labeled classical.

The second is the labels. Every song has only one label, and songs that were listed under more than one genre were removed. A model can fit its labels well and still miss what it is meant to do, like the health algorithm that looked at cost instead of real need. \enquote{The Champion} by Carrie Underwood and Ludacris is labeled country, but the model is almost 97\% sure it is hip hop, and to be fair, Ludacris does rap on it.

The third is proxies. Popularity was left out on purpose, because it says something about fame and not about sound, even though the accuracy dropped from 88\% to 82\%. A model can still pick up the same kind of thing through other features, and \enquote{explicit} is one of them: a quarter of the hip hop songs are explicit, but none of the classical songs are. It is also the feature the model uses the least, so it could simply be removed.

The fourth is feedback loops. If the labels were used to build playlists, a song shown to the wrong listeners would get fewer plays, look less popular and then be shown even less.

The fifth is validation and monitoring. Most of the test songs are by artists the model has already seen in training. For new artists it is right only about 70\% of the time, compared to 81\% for the others, so the results look better than they would in real use. The sepsis model from the lecture, which did worse in hospitals than reported, is an example of the same problem. On top of that, genres change over time, and since November 2024 Spotify no longer gives these audio features to new apps, so the model would be hard to keep up to date.

\section{Choosing the thresholds}

The usual threshold of 0.5 assumes that the classes are balanced and that every mistake costs the same, and neither is true here, since there are 991 heavy metal songs but only 500 jazz songs. That is why every genre gets its own threshold, which is the lowest probability it needs before the app picks it, and if no genre gets there, the app says \enquote{not sure}. To find the thresholds, values from 0.05 to 0.95 were tested with cross validation on the training songs, and the one with the best F1 score was kept for each genre. Country ended up at 0.30, classical and jazz at 0.35, techno at 0.40, hip hop at 0.45 and heavy metal at 0.55. Like the lecture recommends, the thresholds are stored separately from the model, so they can be changed without training it again.

A good example is \enquote{justified} by Kacey Musgraves. It is labeled country, but the model gives it about 30\% hip hop and 17\% country. Since no genre reaches its threshold, the app says \enquote{not sure} instead of giving a wrong answer, and that happened for 37 of the 930 test songs. When the app does give an answer, it is right 83\% of the time, compared to 82\% without thresholds, but if you count \enquote{not sure} as wrong, the accuracy goes down to 80\%. So the thresholds make the model more careful, not smarter.

Picking a threshold is really a choice about who pays for the mistakes. A low threshold finds more songs but also gives more wrong labels, and a high threshold does the opposite. Because country has a low threshold, the app calls more songs country, and more of those answers are wrong. Heavy metal has a high threshold, so the app is right more often when it says heavy metal, but it misses more of the real heavy metal songs. F1 treats both kinds of mistakes as equally bad, so if one of them mattered more, a score like F beta, or tuning based on what each mistake costs, would fit better. \enquote{Not sure} is not spread evenly either, since it hits about 7\% of the hip hop songs but less than 1\% of the classical ones. Thresholds can also make the results look better than they are. With every threshold at 0.70, the app is right 92\% of the times it answers, but it says \enquote{not sure} for 245 of the 930 songs. That is why the app counts \enquote{not sure} as a wrong answer and always shows how many songs got it.

\section{Is it fair, and can it be explained?}

The lecture shows that thresholds for each group can be tuned to meet a fairness criterion, but tuning for F1 does not meet any of them. The app finds 89\% of the classical songs but only 61\% of the jazz songs, so there is no equal opportunity, and how often it is right when it names a genre also varies a lot between the genres. Inside classical, the gap between Western and Indian songs is even bigger. Since the fairness criteria cannot all be met at the same time, one of them has to be chosen, and for artists equal opportunity seems the fairest, because a song that gets missed hurts the artist behind it.

Threshold tuning also assumes that the probabilities can be trusted. The model was never calibrated, for example with Platt scaling, but its confidence still matches how often it is right quite well. When it is at least 90\% sure, it is right 97\% of the time, and when it is less than 50\% sure, only 44\% of the time.

XGBoost is a black box, so it needs extra tools to be explained. Permutation importance and the total gain from XGBoost both say that acousticness, danceability and instrumentalness matter the most, although features that go together, like energy and loudness, can end up sharing the importance. This kind of overall explanation helps the people who build and check the model, but an artist would want to know why their song got the genre it did. SHAP or LIME could show that, but they explain an approximation of the model and not the model itself. The lecture also points out that GDPR and the AI Act ask for explanations, which is one more reason not to use the app for decisions about people.

\section{Can the app be trusted?}

Several things make the app easier to trust. It shows its own test results, puts the real label next to each prediction, warns when a song was used in training, says \enquote{not sure} instead of guessing and rejects wrong input. The results can also be repeated, since the random state is fixed, the model files are saved and the code is in git. Other things work against trust. There is no login, so if one person changes a threshold, it changes for everyone. Song lookups also go to outside services like Spotify and YouTube, which matters for privacy under GDPR, even though the model itself uses no personal data. The polished design can also make the model seem more reliable than it is. There is no plan for monitoring the model yet, and if it were ever used for decisions about artists, or the gaps between groups grew, that would be a reason to stop. All in all, it can be trusted as a hint for clear songs when it is very sure, but not for mixed music, music from outside the West, new artists or answers below 70\%.

\section{How easy is it to use?}

The app has three views: Predict, Thresholds and Features. You can search for a song or describe one with sliders, and you see the prediction next to the real label, bars that show how sure the model is with a line for each threshold, and a short explanation when the answer is \enquote{not sure}. The album covers, videos and DJ mode make it fun to use, but they can also take attention away from the model's limits, and the light grey text can be hard to read. It also needs two servers, Python 3.12 or newer and Node.js 20 or newer to run, so it is more of a demo than a finished product.

\section{Conclusion}

The stakes are low and the app is open about its results, but the model is weak, and the overall accuracy hides that most of the mistakes land on Indian music, songs that mix genres and new artists. The thresholds make it more careful, but every threshold also decides who pays for the mistakes. A better version would need fairer data, testing on new artists, a clear fairness criterion like equal opportunity, explanations for single songs and a login with proper monitoring.

\section{Sources}

{\footnotesize
\begin{itemize}[itemsep=1pt,topsep=2pt]
  \item Lecture slides: \enquote{Production Use, MLOps, Ethics \& Interpretability}, ITF31519, {\O}stfold University College.
  \item Spotify Tracks Dataset (Kaggle): \url{https://www.kaggle.com/datasets/maharshipandya/-spotify-tracks-dataset}
  \item NIST AI Risk Management Framework: \url{https://doi.org/10.6028/NIST.AI.100-1}
  \item The EU AI Act (European Commission): \url{https://digital-strategy.ec.europa.eu/en/policies/regulatory-framework-ai}
  \item EU data protection and GDPR (European Commission): \url{https://commission.europa.eu/law/law-topic/data-protection/data-protection-eu_en}
  \item Gender Shades (Buolamwini and Gebru, 2018): \url{https://proceedings.mlr.press/v81/buolamwini18a.html}
  \item Racial bias in a health algorithm (Obermeyer et al., 2019): \url{https://doi.org/10.1126/science.aax2342}
  \item The Epic sepsis model (Wong et al., 2021): \url{https://doi.org/10.1001/jamainternmed.2021.2626}
  \item Why fairness criteria cannot all hold at once (Kleinberg et al., 2017): \url{https://arxiv.org/abs/1609.05807}
  \item Spotify Web API changes (November 2024): \url{https://developer.spotify.com/blog/2024-11-27-changes-to-the-web-api}
  \item Probability calibration (sklearn user guide): \url{https://scikit-learn.org/stable/modules/calibration.html}
\end{itemize}}

\end{document}
